\documentclass[12pt]{article}

\usepackage[a4paper, total={6in, 8in}]{geometry}
\usepackage{textcomp}

\begin{document}
\setlength{\parindent}{0pt}

\def \t {\textrightarrow}
\def \v {\vspace{1em}}
\def \bi {\begin{itemize}}
\def \ei {\end{itemize}}
\def \s[#1] {\section*{#1}}
\def \ss[#1] {\subsection*{#1}}
\def \sss[#1] {\subsubsection*{#1}}

\s[Recuperare 16 marzo]
\s[The wide sargasso sea - Jean Rhys]
She is a postcolonial author \t and writes "The wide sargasso sea", a postcolonial novel \t west indies is the carabbieans \\
Antoinette had a tragic childhood \t her father died, she was neglected by her mother and had financial problems \t but starting from 1963, the economic situation became worse because slavery was abolished \t luckly she married a rich man \\
Bertha never speaks in the novel \t she is only a crazy and animalesque woman (in Jane Eyre) \t here she is beatufil \\
No onmiscent narrators \t and the novel has 3 parts \t Rochester is a vile goldigger and is not a reliable as a narrator, and Bertha too (she has mental issues) \\
Antoinette had burnt the Thornfield Hall \\
It is a prequiel and a rewriting of Jane Eyre \t part three is a rewriting, the facts of the first 2 parts take place before

\v

Her mother was a Creole, term thas has 3 meanings:
\bi
  \item white person of european descent, however born and raised in the colony \t Antoinette is an example
  \item it can also be of mixed ethnicity
  \item creole languages are a blend of one or even two european languages and a native one
\ei
Jean Rhys was too progressive and was not appreciated \\
Rochester felt a sense of displancement as he was in the West Indies \t while Antoinette felt displaced in England

\ss[A disappointing colonial marriage - extract]
In part two the narrator is Rochester \\
Rocherster is describing his newly married wife (antoinette) and describes how she talks to Caroline, a servant \t he is not used to being this nice to servants \\
He imagines talking to his father \t he says that he has accomplished something: he has married a rich girl \t Rocherster was a disappointement to his father \\
Constant rain that is not puring = drizzle \t but here the rain is intense \\
His british mindset finds inacceptible the close relationship between Antoinette and her servant \\
He also disapproves of Antoinette to be too wild \t she runs in the rain to meet the servant \\
He describes her eyes as alien = different \\
In Jane Eyre Bertha is of mixed race \t so there are racist remarks by some characters \t here Antoinette is white, but only had european origins \t she is still inferior for this reason \\
Antoine brought to the table 30.000 pounds as the dowry \t he sees the marriage as a financial transaction \t he marries her for the money \\
He despise both creoles and natives \t Mr. Rocherster here is voulnerable, has no moral \t he is not presented like an imperialistic hero

\ss[Extract from part III]
Here Antoinette is the narrator \\
She is an alcoholic and abuses Antoinette \\
Sometines Antoinette has visions \t and when she was a kid, she sees her mother \\
When she went to England, she was locked in the attic \t she thought that London was an illusion
\end{document}
